\begin{lemma}[Claim 4.5, singleton support-class branch]
Let \(G\) be a finite loopless multigraph, let \(e_0\) be the
distinguished edge, and let
\[
 T_j(v_n)\preccurlyeq T_j\preccurlyeq T_n,\qquad
 T_j(v_n)\preccurlyeq T'_j
\]
be tree-sequences.  The last displayed vertex of \(T_j(v_n)\) is
\(v_n\), and \(v_n\) belongs to both \(T_j(v_n)\) and \(T_j\).
Let \(\varphi_{j-1}\) and \(\varrho_1\) be partial \(k\)-edge-colorings,
let \(D_{j-1}\) and \(S'\) be finite color sets, and suppose that
\(\varrho_1\) is proper and leaves exactly \(e_0\) uncolored.  Assume
that \(T'_j\) is a \(\varrho_1\)-\noderef{TashkinovClosure} of
\(T_j(v_n)\) and that \(V(T'_j)\) is
\(\varrho_1\)-\noderef{ColoredVertexSetElementary}.

Let \(\gamma_{1_1},\gamma_n,\delta_n,\beta\) be colors satisfying
\[
 \gamma_{1_1}=\gamma_n,\qquad
 \gamma_n\ne\delta_n,\qquad
 \beta\ne\gamma_n,\qquad \beta\ne\delta_n.                    \tag{1}
\]
The first equality is the color-chain identity for the singleton class
\(L_1=\{n\}\); in particular, no preterminal slot is assumed.  Let
\(u\in V(T'_j)\setminus V(T_n)\), with
\[
 \beta\in\overline{\varrho}_1(u),\qquad
 \delta_n\in\overline{\varrho}_1(v_n).                         \tag{2}
\]
Let \(P\) be a specified simple walk in the
\(\noderef{TwoColorIncidenceGraph}\) of
\((\gamma_n,\delta_n)\), from \(v_n\) to \(x\), whose support is exactly
the \(\noderef{TwoColorComponent}\) through \(v_n\), and assume
\(x\notin V(T_n)\cup V(T'_j)\).

Retain the following exact exit-suffix data.  There is a specified simple
walk \(Q\) from \(x\) to \(u\), of length at least two, whose support is
contained in that of \(P\); for every
\(w\in V(T_n)\cup V(T'_j)\),
\[
 \mathsf{inl}(w)\in\operatorname{supp}(Q)\quad\Longleftrightarrow\quad
 w=u.                                                         \tag{3}
\]
There is also an exact edge identity
\[
 g\in\noderef{BoundaryColorEdges}_{\varrho_1,\gamma_n}
       (V(T_n)\cup V(T'_j))
\]
which is incident with \(u\) and whose edge-node lies on \(Q\).

Assume the Claim 4.3 transport data
\[
 \varrho_1\text{ is }
 (T_j(v_n)-v_n,D_{j-1},\varphi_{j-1})
 \text{-}\noderef{StableColoringBeforeVertex},                 \tag{4}
\]
\[
 \overline{\varrho}_1(w)=\overline{\varphi}_{j-1}(w)
 \quad(w\in V(T_j)\setminus\{v_n\}),                           \tag{5}
\]
and
\[
 \varrho_1(T_j.\operatorname{edge}(a))
 =\varphi_{j-1}(T_j.\operatorname{edge}(a))
 \quad(|T_j(v_n)|\le a).                                      \tag{6}
\]
Suppose \(T_j\) is obtained from \(T_j(v_n)\) by a
\(\varphi_{j-1}\)-\noderef{TashkinovAugmentationSequence}, and
\[
 \delta_n\notin\noderef{TreeMissingColors}(\varphi_{j-1},T_j). \tag{7}
\]

Finally, let \(z\in V(T_j)\) occur at the exact displayed position
\(q_z\ge |T_j(v_n)|\).  Assume the usual Observation 23 alternative:
either \(q_z>|T_j(v_n)|\), \(\beta\) is missing at \(z\), and no strictly
intermediate displayed vertex misses \(\beta\), or \(q_z\) is the last
position and no later displayed vertex misses \(\beta\).  In either case
assume the full exact exclusions
\[
 \beta\notin\overline{\varrho}_1(w)
 \quad(w\in\noderef{TreeVerticesBefore}(T_j,z)),               \tag{8}
\]
and
\[
 \varrho_1(T_j.\operatorname{edge}(a))\ne\beta
 \quad(|T_j(v_n)|\le a<q_z).                                  \tag{9}
\]

Define
\[
 \varrho_0=\noderef{OutsideTreeColorSwap}
   (T'_j,\varrho_1,\beta,\delta_n),\qquad
 \varrho_3=\noderef{TwoColorComponentSwap}
   (G,\varrho_0,\gamma_n,\beta,\mathsf{inl}(u)).               \tag{10}
\]
Then \(\varrho_3\) is proper, leaves exactly \(e_0\) uncolored, and has
the following properties.
\begin{enumerate}
\item[(d1)] It is
\((T_j(v_n)-v_n,D_{j-1},\varphi_{j-1})\)-stable.
\item[(d2)] If
\[
 w\in
 \bigl((V(T_j(z))\cup V(T'_j(u)))\setminus\{u\}\bigr)
       \setminus\{z\},
\]
where the two displayed prefixes are represented by their exact
before-vertex sets together with their endpoints, then
\[
 \overline{\varrho}_3(w)=\overline{\varrho}_1(w).
\]
Every exact edge \(T_j.\operatorname{edge}(a)\) with
\(|T_j(v_n)|\le a<q_z\), and every exact edge
\(T'_j.\operatorname{edge}(a)\) lying strictly before an occurrence of
\(u\), has the same color under \(\varrho_3\) and \(\varrho_1\).
Moreover,
\[
 \beta\in\overline{\varrho}_1(z)
 \Longrightarrow \delta_n\in\overline{\varrho}_3(z),          \tag{11}
\]
and
\[
 z\in\noderef{TreeVerticesBefore}(T'_j,u)
 \Longrightarrow
 \overline{\varrho}_3(z)=\overline{\varrho}_1(z).              \tag{12}
\]
\item[(d3)] \(\gamma_{1_1}\in\overline{\varrho}_3(u)\).
\end{enumerate}
All assertions retain exact edge identities, including identities of
parallel edges.
\end{lemma}
\begin{proof}
The prefix relations, expanded using
\noderef{TreeSequenceInitialSegment} and
\noderef{TreeSequenceVertexSet}, put \(v_n\) in \(V(T'_j)\) and put
\(V(T_j)\) inside \(V(T_n)\).  Thus (2) and the union definition
\noderef{TreeMissingColors} make both \(\beta\) and \(\delta_n\)
tree-missing on \(T'_j\).  The closure clause in
\noderef{TashkinovClosure}, interpreted by
\noderef{ColoredVertexSetClosed}, says that these colors have no boundary
edge at \(V(T'_j)\).  Since they are distinct,
\noderef{OutsideTreeColorSwapStable} applies to (10).  Consequently
\(\varrho_0\) is proper, leaves exactly \(e_0\) uncolored, is stable on
\(T'_j\) with control set \(S'\), and satisfies
\[
 \varrho_0(e)=\varrho_1(e)
 \quad\text{for every exact edge identity incident with }V(T'_j). \tag{13}
\]

We identify the component \(K\) switched in the definition of
\(\varrho_3\).  Every edge-node of \(Q\) has
\(\varrho_1\)-color \(\gamma_n\) or \(\delta_n\), because \(Q\) is
contained in the support of the full old component \(P\), in the sense
of \noderef{TwoColorComponent} and
\noderef{TwoColorIncidenceGraph}.  If a \(\delta_n\)-edge-node of \(Q\)
were incident with \(T'_j\), its other incident vertex-node would also
lie on \(Q\).  By (3) that vertex must be \(u\); simplicity of \(Q\)
then forces this to be its terminal edge identity.  But the specified
terminal identity \(g\) has color \(\gamma_n\), by the definition
\noderef{BoundaryColorEdges}, and properness
\noderef{IsProperEdgeColoring} excludes a second
\(\gamma_n\)-identity at \(u\).  Therefore (13) fixes the terminal
identity, while every \(\delta_n\)-identity on the remaining suffix lies
outside \(T'_j\) and is changed to \(\beta\) by
\noderef{OutsideTreeColorSwap}; all \(\gamma_n\)-identities remain fixed.
It follows edge by exact edge that \(Q\) is a
\((\gamma_n,\beta)\)-walk under \(\varrho_0\).

We now prove that its support is the whole new component:
\[
 \operatorname{supp}(Q)
 =\noderef{TwoColorComponent}
   (G,\varrho_0,\gamma_n,\beta,\mathsf{inl}(u)).               \tag{14}
\]
One inclusion follows by connectivity of the walk just proved.  For the
reverse inclusion, consider possible new adjacencies along \(Q\).
At \(x\), every new \(\gamma_n\)- or \(\beta\)-adjacency yields an old
\(\gamma_n\)- or \(\delta_n\)-adjacency: an edge meeting \(T'_j\) is
fixed by (13), while a boundary \(\beta\)-edge is forbidden by closure
and the fact that \(\beta\) is missing at \(u\).  Since \(P\) represents
the full old component and is simple with endpoint \(x\), this adjacency
is the unique first edge-node of \(Q\).  At \(u\), \(\beta\) is missing
under \(\varrho_0\), and properness makes the
\(\gamma_n\)-colored identity \(g\) its unique new two-color neighbor.
Each internal node of the simple path \(Q\) already has two distinct
neighbors.  The theorem
\noderef{TwoColorIncidenceDegreeAtMostTwo} bounds every degree in this
incidence graph by two because \(\varrho_0\) is proper.  Hence no
internal node has a third neighbor and neither endpoint has an additional
neighbor.  The support of \(Q\) is therefore closed under adjacency.
Every node connected to \(u\) lies in it, proving (14).

Combining (14) with (3) gives the exact intersection
\[
 K\cap\{\mathsf{inl}(w):w\in V(T_n)\cup V(T'_j)\}
 =\{\mathsf{inl}(u)\}.                                       \tag{15}
\]
If an exact edge identity in \(K\) were incident with a vertex
\(w\in V(T_n)\cup V(T'_j)\), adjacency in
\noderef{TwoColorIncidenceGraph} would put \(\mathsf{inl}(w)\) in the
same \noderef{TwoColorComponent}.  Thus (15) says that the second switch
fixes every edge incident with this union away from \(u\).
Since \(\gamma_n\ne\beta\), \noderef{TwoColorComponentSwapProper} proves
that \(\varrho_3\) is proper.  Expanding
\noderef{TwoColorComponentSwap} shows that it merely interchanges two
nonempty option values on \(K\) and fixes all other values; together with
\noderef{IsOneEdgeUncolored}, this proves that exactly \(e_0\) remains
uncolored.

Every vertex strictly before \(v_n\) in \(T_j\) belongs to the seed:
the seed ends at \(v_n\), its inclusion in \(T_j\) is exact, and displayed
vertices are injective.  This follows directly from
\noderef{TreeSequenceInitialSegment},
\noderef{TreeSequenceVertexSet}, and the before relation
\noderef{TreeVerticesBefore}.  Such a vertex lies in \(V(T'_j)\), is
different from \(u\), and hence every incident identity is fixed by both
swaps, using (13) and (15).  The witness definition
\noderef{PresentColors} and its complement
\noderef{MissingColors} give equality of its old and new missing-color
sets.  Both the incident-edge and missing-set clauses of (4) therefore
transfer from \(\varrho_1\) to \(\varrho_3\), proving (d1).

We next prove the exact edge clause in (d2).  Induct along the finite
chain \noderef{TashkinovAugmentationSequence}.  The one-step statement
\noderef{TashkinovAugmentation} says that each exact displayed edge
appended after the seed has, under \(\varphi_{j-1}\), a color missing at
an earlier displayed vertex.  By (7) that color is not \(\delta_n\);
by (6) its \(\varrho_1\)-color is the same, and (9) separately excludes
\(\beta\) before \(q_z\).  Hence the outside transposition fixes each
such exact identity.  Seed identities occur in \(T'_j\), so (13) fixes
them too.  Their endpoints lie in \(V(T_j)\subseteq V(T_n)\), by the
attachment law of \noderef{TreeSequence}; (15) therefore makes the second
switch fix them.  This proves the \(T_j\)-edge clause.  An exact displayed
edge of \(T'_j\) has two displayed endpoints by that same attachment law.
Equation (13) fixes it under the first switch.  If it lay in \(K\), both
endpoint vertex-nodes would lie in \(K\); (15) permits only \(u\), while
looplessness forbids both endpoints being \(u\).  Hence the second switch
also fixes every required \(T'_j\)-edge identity.

For the missing-set clause, first take a listed vertex in
\(V(T'_j)\setminus\{u\}\).  Equations (13) and (15) fix every incident
identity, so \noderef{PresentColors} and \noderef{MissingColors} give the
required equality.  Now take a listed vertex \(w\) of \(T_j(z)\) outside
\(T'_j\), with \(w\ne z\).  By (8), \(\beta\) is present at \(w\).
Also \(w\ne v_n\), because \(v_n\in V(T'_j)\); (7) says that
\(\delta_n\) is present under \(\varphi_{j-1}\), and (5) transfers that
presence to \(\varrho_1\).  Both colors are tree-missing on the closed
set \(V(T'_j)\), so no incident edge of either color can join \(w\) to
that set, by \noderef{BoundaryEdges} and
\noderef{ColoredVertexSetClosed}.  The outside swap therefore transposes
the two complete, nonempty incidence sets at \(w\) and fixes all other
colors.  Its present set, hence its missing set, is unchanged.  Finally
\(w\in V(T_n)\) and \(w\ne u\), so (15) says the component swap fixes all
its incident identities.  This proves the vertex equality in (d2).
The same incident-edge argument immediately proves (12) when
\(z\) is displayed strictly before \(u\) in \(T'_j\).

Suppose now that \(\beta\in\overline{\varrho}_1(z)\).  Since
\(z\in V(T_j)\subseteq V(T_n)\) but \(u\notin V(T_n)\), the vertices are
distinct.  Elementariness of \(V(T'_j)\), together with the simultaneous
missingness of \(\beta\) at \(u\), implies \(z\notin V(T'_j)\).  Hence
\(z\ne v_n\); equations (7) and (5) show that \(\delta_n\) is present at
\(z\).  The same closed-boundary argument shows that all incident
\(\beta\)- and \(\delta_n\)-identities lie outside \(T'_j\).  Thus
\noderef{OutsideTreeColorSwap}, together with
\noderef{PresentColors} and \noderef{MissingColors}, makes
\(\delta_n\) missing at \(z\).  Equation (15) makes the second switch
fix every identity incident with \(z\), proving (11).

Finally, (13) preserves the missingness of \(\beta\) at \(u\) and the
\(\gamma_n\)-colored identity \(g\).  In the rooted component switch,
every incident \(\gamma_n\)-identity belongs to \(K\) and becomes
\(\beta\), while no incident \(\beta\)-identity exists beforehand.
Expansion of \noderef{TwoColorComponentSwap},
\noderef{PresentColors}, and \noderef{MissingColors} therefore gives
\(\gamma_n\in\overline{\varrho}_3(u)\).  The singleton identity
\(\gamma_{1_1}=\gamma_n\) from (1) is exactly (d3).
\end{proof}

